\begin{abstract}
\noindent    
I describe the development and evaluation of an open-source tool for detecting hallucinated citations in academic papers, built entirely on free public APIs (CrossRef, OpenAlex, Semantic Scholar). Testing against 391 real citations and 500 fabricated ones, I find that deterministic validation achieves 0\% false-positive rate and 100\% detection on citations with DOIs, but is nearly blind to fabricated citations without DOIs (4\% detection). Adding AI analysis (Gemini 2.5 Flash, free tier) raises detection to 94\% on fakes, but introduces a 72\% false-positive rate on real arXiv preprints.

The central finding is not about the tool. It is about an epistemological error that recurs at every tier of detection: the conflation of unverifiable with fabricated. I identified and fixed this error in my deterministic logic, only to discover that AI models reproduce it independently. Grounded AI's analysis for Nature shows the same pattern — a 22\% false-flag rate on their most suspicious publications.

I argue that citation hallucination detection is fundamentally a problem of epistemology, not engineering, and that any system, whether rule-based, AI-powered, or human, must grapple with the distinction between absence and fabrication. The tool, code, data, and all experimental results are freely available.
\\\\   
\smallskip
\noindent\textbf{Keywords:} citation hallucination, generative AI, open-source tools, scientific integrity, epistemology, BibTeX validation
\end{abstract}